\begin{frame}{Diffie-Hellman-Merkle-Schlüsseltausch (DHM)}
	1976 erfanden Whitfield Diffie, Martin Hellman und Ralph Merkle das erste Verfahren für einen sicheren Schlüsselaustausch über unsichere Kanäle.
	
	\begin{myidea}[Diskretes Logarithmusproblem]
		Das Verfahren basiert auf einer mathematischen \textbf{Einwegfunktion} modulo einer Primzahl $p$:
		\begin{itemize}[<+->]
			\item \textbf{Einfach}: $x = g^a \bmod p$ berechnen.
			\item \textbf{Sehr schwer}: Aus $x, g, p$ den geheimen Exponenten $a$ zurückrechnen (\emph{Diskreter Logarithmus}).
		\end{itemize}
	\end{myidea}
\end{frame}

\begin{frame}{Ablauf des DHM-Protokolls}
	\begin{myprocedure}[Protokoll \acrshort{DHM}]
		{\small
		Öffentlich vereinbarte Werte: Primzahl $p$ und Basis $g < p$.
		
		\begin{enumerate}[<+->]
			\item \textbf{Alice}: Wählt geheime Zahl $a < p$. Berechnet $x := g^a \bmod p$ und sendet $x$ an Bob.
			\item \textbf{Bob}: Wählt geheime Zahl $b < p$. Berechnet $y := g^b \bmod p$ und sendet $y$ an Alice.
			\item \textbf{Alice} berechnet gemeinsamen Schlüssel:
			      $s_{AB} := y^a \bmod p$.
			\item \textbf{Bob} berechnet gemeinsamen Schlüssel:
			      $s_{AB} := x^b \bmod p$.
		\end{enumerate}
		}
	\end{myprocedure}
\end{frame}

\begin{frame}{Warum erhalten beide denselben Schlüssel?}
	Aus den Potenzgesetzen folgt:
	\[ (g^b)^a \bmod p = g^{b \cdot a} \bmod p = g^{a \cdot b} \bmod p = (g^a)^b \bmod p \]
	
	\pause
	\begin{align*}
		\text{Alice berechnet: } & y^a \bmod p = (g^b)^a \bmod p = \mathbf{g^{ab} \bmod p} \\
		\text{Bob berechnet: }   & x^b \bmod p = (g^a)^b \bmod p = \mathbf{g^{ab} \bmod p}
	\end{align*}
	
	\pause
	\begin{alertblock}{Was sieht der Abhörer Eve?}
		Eve sieht nur $p, g, x, y$. Ohne $a$ oder $b$ kann Eve $g^{ab} \bmod p$ nicht berechnen!
	\end{alertblock}
\end{frame}

\begin{frame}{Rechenbeispiel: DHM-Schlüsseltausch}
	\begin{exampleblock}{Beispiel mit kleinen Zahlen}
		Öffentliche Parameter: $p = 13$, $g = 2$.
		
		\begin{enumerate}[<+->]
			\item \textbf{Alice} ($a = 5$):
			      $x = 2^5 \bmod 13 = 32 \bmod 13 = 6$
			\item \textbf{Bob} ($b = 4$):
			      $y = 2^4 \bmod 13 = 16 \bmod 13 = 3$
			\item \textbf{Alice}:
			      $s_{AB} = y^a \bmod p = 3^5 \bmod 13 = 243 \bmod 13 = \mathbf{9}$
			\item \textbf{Bob}:
			      $s_{AB} = x^b \bmod p = 6^4 \bmod 13 = 1296 \bmod 13 = \mathbf{9}$
		\end{enumerate}
	\end{exampleblock}
\end{frame}

\begin{frame}{Auftrag: Diffie-Hellman-Merkle}
	\begin{itemize}
		\item \aufgabebox{Lösen Sie die Aufgaben 3.2, 3.3 und 3.4 im Skript}
		\item Zeit: 15 Minuten
	\end{itemize}
\end{frame}
